
Multi-dimension trustworthiness benchmarks for large language models, such as DecodingTrust \citep{Wangetal2023} and TrustLLM \citep{Sunetal2024}, evaluate failure modes including adversarial robustness, calibration, hallucination, and safety independently, without testing whether performance on one dimension predicts performance on another. An implicit assumption in such frameworks is that these failure modes are positively coupled — that a model fragile under adversarial perturbation will also be miscalibrated and unreliable in other respects.

This study examines whether that assumption holds for the adversarial fragility–calibration pair. A key methodological challenge in such an analysis is confounding by general capability: larger, more capable models tend simultaneously to perform better on adversarial robustness benchmarks and to exhibit lower calibration error. Any raw correlation between adversarial robustness and calibration across a diverse model set is therefore confounded by the underlying capability signal. To address this, the study introduces Residual Instability (RI), defined as the OLS residual of AdvGLUE accuracy drop after regressing out a PCA-derived capability index (PC1) and mean model confidence. RI is constructed to be orthogonal to general capability by design.

The primary empirical finding is that RI significantly anticorrelates with Expected Calibration Error (ECE) on the arc_challenge benchmark: ρ = −0.535, p = 0.0034, 95\% CI = [−0.782, −0.101], n = 30. This result is in the opposite direction from the coupled failure cascade prediction, which expected a positive partial correlation of ρ ≥ +0.4. The finding indicates that, after controlling for capability, models with higher adversarial fragility tend to be better calibrated, not worse.

The study makes three contributions: (1) the RI construct — a reusable, capability-orthogonal measure of adversarial fragility (R² = 0.529, SD = 0.121, VIF = 1.000); (2) the empirical RI–ECE anticorrelation across 30 LLMs from 9 families; and (3) a calibration–robustness trade-off hypothesis as a testable alternative mechanistic explanation grounded in the RLHF literature.

The study scope is explicitly limited to the RI–ECE dimension. Sub-hypotheses testing whether RI predicts hallucination rate (HaluEval), safety failure (HarmBench), and output variance (OVI-GSM8K) were blocked by the gate failure on the primary mechanism hypothesis and were not executed. This limitation is a central constraint on the conclusions that can be drawn.

The paper is organized as follows. Section 2 reviews related work on multi-dimension trustworthiness evaluation, adversarial robustness, and model calibration. Section 3 describes the RI methodology. Section 4 details the experimental setup. Section 5 presents results. Section 6 discusses mechanistic interpretations and limitations. Section 7 concludes.

